\chapter{النماذج}
\label{ch:builder-forms}

\section{ماذا ستبني في هذا الفصل؟}
\label{sec:builder-forms-intro}

النموذج صفحة أسئلة يعبّئها الموظف أو مقدّم الطلب، وتصل إجاباته إلى من يراجعها، وقد تُكتب في نوع سجلات. ويشرح هذا الفصل بناء النموذج من أوله إلى متابعة طلباته: إنشاؤه، وترتيب صفحاته وأسئلته، وضبط قواعده وأزراره، وربطه بالسجلات، وتحديد من يصل إلى ردوده، ثم رسم مسار الطلب وموافقاته، وتجربته ونشره.

ويُحفظ كل تعديل تجريه في «منشئ النماذج» في مسودة، أما الناس فيعبّئون دائماً الإصدار المنشور. لذلك تستطيع أن تعدّل وتختبر بأمان، ولا يصل التغيير إلى أحد إلا بعد النشر.

\begin{rcnote}[مثال: شركة الأفق]
ينشئ مسؤول النظام نموذج «طلب زيارة فرع»، يعبّئه عمر حداد فيُنشأ من إجاباته سجل في «طلبات الزيارة» بعد أن تعتمده سارة الخطيب.
\end{rcnote}

\begin{figure}[H]
\centering
\begin{tikzpicture}[node distance=0.5cm]
\node[stepBox] (a) {\RL{المحتوى: الصفحات والأسئلة}};
\node[stepBox,below=of a] (b) {\RL{السلوك: القواعد والأزرار}};
\node[stepBox,below=of b] (c) {\RL{الربط: السجلات والوصول}};
\node[stepBox,below=of c] (d) {\RL{المسار والموافقات}};
\node[stepBox,below=of d] (e) {\RL{الاختبار والنشر}};
\node[stepBox,below=of e] (f) {\RL{الطلبات والمتابعة}};
\draw[flowsTo] (a) -- (b);
\draw[flowsTo] (b) -- (c);
\draw[flowsTo] (c) -- (d);
\draw[flowsTo] (d) -- (e);
\draw[flowsTo] (e) -- (f);
\end{tikzpicture}
\rclegend
\caption{مراحل بناء النموذج}
\label{fig:builder-forms-overview}
\shotnote{يرتّب المخطط صفحات الفصل بترتيب البناء المقترح: ما يُسأل أولاً، ثم كيف يتصرف النموذج، ثم أين تذهب الإجابات ومن يراها. وتأتي المتابعة أخيراً لأنها لا تبدأ إلا بعد نشر النموذج ووصول الطلبات.}
\end{figure}

\section{صفحات هذا الفصل}
\label{sec:builder-forms-pages}

تبدأ الصفحة الأولى من «النماذج» في قسم «الأساسية»، أما بقية الصفحات فتشرح أدوات «منشئ النماذج» التي تصل إليها بفتح النموذج، والصفحة الأخيرة تشرح جدول الطلبات. واتبع الترتيب في المرة الأولى، ثم عُد إلى أي صفحة عند الحاجة.

\begin{table}[H]
\centering
\caption{بناء النموذج خطوة بخطوة}
\label{tab:builder-pages-forms}
\rcZebra
\begin{tabularx}{\linewidth}{R{4.2cm} Y}
\rcHeaderRow
\thh{الصفحة وقسمها} & \thh{ماذا تقدم هذه الصفحة؟}\\
النماذج — الأساسية & تعرض نماذج المؤسسة، ومنها تنشئ نموذجاً جديداً أو تنسخه أو تستورده.\\
منشئ النماذج وأدواته — منشئ النماذج & لوحة النموذج وشريط «أدوات النموذج»، ومساعد الذكاء الاصطناعي الذي يقترح التعديلات.\\
الصفحات والأقسام والجداول — منشئ النماذج & تقسيم النموذج إلى صفحات، وتجميع الأسئلة في أقسام متكررة أو جداول.\\
أنواع الأسئلة: النص والاختيار والأرقام — منشئ النماذج & أنواع النصوص والخيارات والأرقام والتقييم، وطرق عرضها.\\
أنواع الأسئلة: التاريخ والموقع والملفات وغيرها — منشئ النماذج & أسئلة الوقت والخرائط والمرفقات والوسائط، واختيار سجل أو طلب موجود.\\
إعدادات السؤال والتحقق — منشئ النماذج & لوحة إعدادات السؤال، وقواعد التحقق ومستوياتها ورسائلها.\\
القواعد — منشئ النماذج & كل ما يُظهر الأسئلة أو يخفيها أو يلزم بها أو يقفلها أو يحسبها، في قائمة واحدة.\\
أزرار الإجراء: عند النقر — منشئ النماذج & أزرار تتحقق وتحفظ وترسل الإشعارات والمستندات بخطوات مرتبة.\\
السجلات: ربط النموذج بنوع سجلات — منشئ النماذج & قراءة الإجابات من السجلات وكتابتها إليها، ومتى تُكتب.\\
الوصول والمشاركون — منشئ النماذج & من يصل إلى كل رد، وماذا يستطيع أن يفعل به، ومن يضغط كل زر.\\
مسار الطلب: الحالات والموافقات — منشئ النماذج & حالات الطلب ومراحله وانتقالاته، وخطوات الموافقة ومراجعة الذكاء الاصطناعي.\\
الإسناد وصفحة العمل — منشئ النماذج & تكليف أشخاص بتعبئة النموذج، ومتابعة التعبئة والموافقات والمتأخر.\\
إعدادات النموذج — منشئ النماذج & الترقيم ومستند PDF ودورة الحياة والمراحل والتعبئة وسياسة النشر.\\
الاختبار والنشر والإصدارات — منشئ النماذج & تجربة النموذج بصفة عضو في مساحة معزولة، ثم نشره ومتابعة إصداراته.\\
الطلبات — قائمة التنقل & جدول ردود النموذج، ومنه تُفتح الطلبات وتُراجع ويُبدأ طلب جديد.\\
\end{tabularx}
\end{table}
